% !TeX TS-program = xelatex
% 虚构演示简历 — 仅用于 hellojobs 测试数据
\documentclass{resume-photo}
\usepackage{xeCJK}
\setCJKmainfont{Noto Serif CJK SC}
\usepackage{fontspec}
\usepackage{enumitem}
\setlist[itemize]{itemsep=0pt, topsep=1pt, parsep=0pt, partopsep=0pt}

\ResumeName{钱峰}

\begin{document}

\ResumeContacts{
  138-0021-0021,%
  \ResumeUrl{mailto:qianfeng@bupt.edu.cn}{qianfeng@bupt.edu.cn},%
  \textnormal{北京邮电大学 | 网络工程 · 学士 | 2020—2024}%
}

\ResumeTitle

\noindent 爬虫与反爬对抗，熟悉 Scrapy、Playwright 与 TLS 指纹；建设合规采集管线，请求成功率从 71\% 提升至 94\%。注重可观测性与文档沉淀，习惯在评审阶段拆解风险；参与线上复盘与跨团队联调，英语 CET-6，可阅读官方文档与 RFC（演示数据）。

\section{教育背景}
\ResumeItem
{北京邮电大学}
[\textnormal{计算机学院 | 网络工程 | 学士}]
[2020.09—2024.06]

\begin{itemize}
  \item 网络协议与系统安全课程为主。
\end{itemize}


\ResumeItem
{实验室与助教}
[\textnormal{本科生科研 / 课程助教}]
[2022—2025]

\begin{itemize}
  \item 进入校级重点实验室参与课题，负责数据采集、清洗与可视化看板搭建。
  \item 担任《数据结构》《操作系统》课程助教，批改作业 200+ 份，答疑课时 40h+。
  \item 组织期中复习串讲与上机辅导，课程评教得分位于前 10\%（演示）。
  \item 整理课程 FAQ 与实验踩坑文档，被后续年级沿用。
\end{itemize}



\section{专业技能}
\begin{itemize}
  \item 熟悉 Python、Scrapy、Redis 去重与分布式调度。
  \item 掌握 Playwright、代理池轮换与验证码打码平台对接（合规场景）。
  \item 了解 JA3 指纹与 HTTP/2 帧级调试。
  \item 熟悉 Linux 日常运维与 Shell，具备日志检索与 CPU/内存突增初步定位能力。
  \item 掌握 Git / CI 与 Code Review，了解 Docker 与 Kubernetes 基本编排与滚动发布。
  \item 了解 OAuth2 / JWT、接口幂等与 REST 规范，具备单元测试与集成测试习惯（JUnit/pytest 等）。
  \item 能阅读英文技术文档，具备跨团队沟通与需求澄清能力（测试简历）。
\end{itemize}

\section{实习经历}
\ResumeItem
{某数据服务商}
[\textnormal{数据采集实习生}]
[2024.01—2024.08]

\begin{itemize}
  \item \textbf{职责}: 电商价格监控链路稳定性。
  \item \textbf{成果}: 日级任务失败重跑率下降 52\%。
\end{itemize}


\ResumeItem
{SHEIN | 全球技术中心}
[\textnormal{供应链后端实习生}]
[2024.02—2024.07]

\begin{itemize}
  \item \textbf{场景}: 跨境库存与仓配调度接口高并发读写，存在热点行更新争用。
  \item \textbf{方案}: 引入分段锁与异步合并写，结合 Binlog 订阅做最终一致。
  \item \textbf{结果}: 热点 SKU 更新 TPS 上限提升 2.4 倍，锁等待时间下降 70\%。
  \item \textbf{工具}: 编写压测脚本与数据构造工具链，纳入 CI nightly。
  \item \textbf{协作}: 与仓储 WMS 团队对齐字段口径，消除跨系统对账差异。
\end{itemize}



\section{项目经历}
\ResumeItem
{合规 Robots 解析器}
[\textnormal{毕业设计}]
[2023.10—2024.05]

\begin{itemize}
  \item 支持多级 robots.txt 缓存与 TTL，误判拦截率 <0.1\%。
\end{itemize}


\ResumeItem
{多租户 SaaS 计费引擎}
[\textnormal{订阅制 + 用量计费（演示）}]
[2023.09—2024.02]

\begin{itemize}
  \item \textbf{模型}: 设计套餐、用量累计、账单周期与 proration 规则。
  \item \textbf{一致}: 使用 Outbox + 消息表保证计费事件与财务对账一致。
  \item \textbf{性能}: 批价任务分片并行，账单生成时间缩短 48\%。
  \item \textbf{测试}: 构造边界用例库 150+，覆盖闰月与升级降级场景。
\end{itemize}



\section{个人荣誉}
\begin{itemize}
  \item 北邮 优秀毕业生
  \item 全国大学生信息安全竞赛 华北三等奖
  \item 全国计算机等级考试四级（网络工程师）（演示）
  \item 校级「优秀学生干部 / 优秀团员」或技术社团骨干（综合测评前 15\%）
  \item 开源贡献：向知名项目提交被合并的 PR（文档/测试/feature）（演示）
\end{itemize}

\end{document}
